所有数值均来自本次附件实跑。基线保留原ALNS与原解码器，统一更正地形模块，运行800次迭代、种子20260924、固定Python散列种子0。改进运行四组偏好，每组两个种子、每个500次迭代，并以基线路线热启动，随后求解有限时长的整数规划。因此下表证明本算例取得了更好的可行方案，不是相同计算预算下的算法排名。
\begin{center}
\begin{tabular}{lrrrrr}\toprule
方案 & 架次 & 完工/min & 能耗/kWh & 加权迟到 & 按期箱数\\\midrule
原ALNS统一地形 & 20 & 159.72 & 64.5185 & 6.6369 & 77/80 \\
改进均衡方案 & 20 & 137.48 & 61.8969 & 0.0000 & 80/80 \\
减少架次方案 & 17 & 155.22 & 60.6636 & 7.2047 & 75/80 \\
\bottomrule\end{tabular}\end{center}
均衡方案比基线缩短完工时间22.24分钟（13.92\%），减少能耗2.6216kWh（4.06\%），架次数相同，全部80箱在期望时间内交付。均衡、完工时间和能耗三组权重最终选择了同一方案。减少架次方案使用17架次，5箱普通物资迟于期望时刻，但全部医疗和首批箱仍满足硬时限。

均衡方案的31个不同硬时限箱（16箱医疗、30箱首批，两组有重叠）全部满足要求。最紧硬时限余量为206.804秒，最低返航SOC为20.971\%，高于20\%下限。每箱恰好一次；具体无人机无任务重叠；具体电池的任务和充电区间无重叠。独立审计使用微秒级端点归一化处理浮点累计误差，不把同一资源前后任务的共同端点计为重叠。
\begin{center}\begin{tabular}{lrrrr}\toprule
机型 & 同时执行无人机峰值 & 库存 & 电池任务加充电峰值 & 电池库存\\\midrule
A & 4 & 4 & 4 & 6 \\
B & 2 & 2 & 4 & 4 \\
C & 2 & 2 & 4 & 4 \\
\bottomrule\end{tabular}\end{center}
独立地形验证将全部120条无向航段与另一套线段矩形相交计算交叉核对，同时检查穿角点、沿网格边界与90\%充电分段边界。修改已导出结果中的能耗、无人机ID或电池ID，校验器能够识别错误。

本次实际增益来自多起点路线搜索及固定路线的机型时序整数规划，后者将均衡ALNS解的完工时间从8909.39秒进一步降至8248.92秒。四次较大的候选池联合模型均在当前时限内未取得可用整数解，因而没有用它们替换已验证方案。程序完整实现该阶段，但不能将“已建立模型”描述为“已取得更优联合解”。固定路线子问题的求解状态如下，均为压紧前的离散模型：
\begin{center}\begin{tabular}{lrrr}\toprule
偏好 & 开始时间列数 & 相对gap & 实际求解/s\\\midrule
均衡 & 3745 & 0.0000\% & 4.02 \\
完工时间 & 4042 & 0.0299\% & 5.80 \\
能耗 & 3130 & 0.3518\% & 1.18 \\
架次 & 2729 & 0.0000\% & 1.57 \\
\bottomrule\end{tabular}\end{center}
均衡固定路线模型的gap为零，但其箱集、顺序和开始网格均受到限制，因此不能称原问题全局最优。最晚返航的统计口径严格为所有运输无人机完成最后一个架次并返回O01的最大时刻。
